\section{Layanan Keamanan Informasi}

Empat masalah Alice--Bob (penyadapan, penyamaran, manipulasi, penyangkalan) dijawab oleh layanan keamanan yang didukung kriptografi \cite{munir_slides_itb,rfc4949,nist_cia}: \textbf{confidentiality}, \textbf{authentication}/\textbf{authenticity}, \textbf{integrity}, dan \textbf{non-repudiation}, dilengkapi \textbf{availability} dalam model CIA. Kasus kebocoran data dan serangan pada bagian urgensi memperkuat motivasi bahwa layanan ini bukan sekadar istilah teoretis, melainkan kebutuhan praktis. Tiga layanan pertama CIA triad dikenal luas; non-repudiation dan authenticity melengkapi model untuk transaksi yang memerlukan kepercayaan dan bukti hukum.

\subsection{Confidentiality (Kerahasiaan)}

Memastikan informasi hanya dapat diakses oleh pihak yang berwenang. Plaintext tidak boleh terbaca oleh Eve meskipun dia memiliki ciphertext. \textbf{Cara mencapai}: enkripsi dengan kunci yang hanya diketahui Alice dan Bob \cite{stallings2017}. Gambar~\ref{fig:layanan-kerahasiaan} menunjukkan plainteks yang menjadi tidak bermakna setelah dienkripsi.

\begin{figure}[htbp]
\centering
\includegraphics[width=0.8\textwidth]{bab-02/kerahasiaan-plain-cipher}
\caption{Kerahasiaan: plainteks (a) dienkripsi menjadi cipherteks (b) yang tidak dapat dipahami pihak yang tidak berhak \cite{munir_slides_itb}.}
\label{fig:layanan-kerahasiaan}
\end{figure}

Confidentiality merupakan fondasi keamanan informasi. Akses ke informasi harus dikontrol untuk mencegah pembagian data tidak sah. Implementasi efektif memerlukan enkripsi data, kontrol akses, dan autentikasi. Ancaman terhadap confidentiality mencakup serangan siber maupun kesalahan manusia seperti berbagi kredensial atau kegagalan mengenkripsi komunikasi.

\subsection{Integrity (Keutuhan)}

Memastikan informasi tidak diubah selama transmisi. Bob harus yakin pesan yang diterima sama persis dengan yang dikirim Alice. \textbf{Cara mencapai}: fungsi hash, Message Authentication Code (MAC), atau tanda tangan digital. Perubahan kecil pun harus terdeteksi, seperti satu digit suhu atau satu bagian foto pada Gambar~\ref{fig:layanan-integritas}.

\begin{figure}[htbp]
\centering
\includegraphics[width=\textwidth]{bab-02/integritas-pesan-diubah}
\caption{Integritas: pesan teks dan citra yang diubah sedikit tampak hampir sama dengan aslinya \cite{munir_slides_itb}.}
\label{fig:layanan-integritas}
\end{figure}

Integrity memastikan data tetap tepercaya dan bebas dari perusakan. Setiap perubahan tidak sah terhadap data harus terdeteksi dan dicegah. Pelanggaran integrity dapat menyebabkan keputusan bisnis yang salah atau kehilangan kepercayaan. Teknologi hash, tanda tangan digital, dan MAC menyediakan mekanisme verifikasi yang diperlukan.

\subsection{Availability (Ketersediaan)}

Memastikan sistem dan data dapat diakses oleh pihak yang berwenang ketika diperlukan. Serangan denial-of-service (DoS) mengancam availability dengan membuat layanan tidak dapat diakses.

Kriptografi secara tidak langsung mendukung availability melalui proteksi terhadap serangan yang menargetkan infrastruktur. Manajemen kunci dan redundansi sistem juga berperan menjaga ketersediaan layanan kriptografi. Dalam praktik, availability sering ditangani melalui kombinasi kebijakan operasional, redundansi jaringan, dan mitigasi DoS.

\subsection{Non-Repudiation (Ketidakbantahan)}

Memastikan pihak yang melakukan tindakan tidak dapat menyangkal bahwa ia telah melakukannya. Alice tidak dapat menyangkal telah mengirim pesan; Bob tidak dapat menyangkal telah menerimanya. \textbf{Cara mencapai}: tanda tangan digital, sertifikat digital, log audit. Gambar~\ref{fig:layanan-nirpenyangkalan} mengilustrasikan nasabah yang menyangkal instruksi transfernya sendiri.

\begin{figure}[htbp]
\centering
\includegraphics[width=0.6\textwidth]{bab-02/nirpenyangkalan-transfer}
\caption{Non-repudiation: tanpa bukti kriptografis, A dapat menyangkal telah meminta transfer ke B \cite{munir_slides_itb}.}
\label{fig:layanan-nirpenyangkalan}
\end{figure}

Non-repudiation penting untuk transaksi hukum, kontrak elektronik, dan pembayaran daring. Tanda tangan digital memberikan bukti matematis yang tidak dapat dipalsukan tentang asal dan integritas pesan. Sertifikat digital dari CA tepercaya mendukung verifikasi identitas dalam skala global \cite{stallings2017}.

\subsection{Authenticity (Keaslian)}

Authenticity memastikan pengirim pesan benar-benar pihak yang mengklaim. Bob yakin pesan berasal dari Alice, bukan Eve yang berpura-pura. \textbf{Cara mencapai}: tanda tangan digital, sertifikat digital, protokol autentikasi.

\begin{figure}[!htb]
\centering
\includegraphics[width=0.5\textwidth]{bab-02/autentikasi-penyamaran}
\caption{Autentikasi: pihak ketiga di tengah mengirim sinyal palsu seolah-olah berasal dari A, sehingga B perlu memverifikasi asal pesan \cite{munir_slides_itb}.}
\label{fig:layanan-autentikasi}
\end{figure}

Gambar~\ref{fig:layanan-autentikasi} menggambarkan masalah ini secara sederhana: B menerima sinyal asap, tetapi tidak dapat memastikan apakah sinyal itu dari A atau dari penyamar.
\begin{table}[htbp]
\centering
\small
\begin{tabular}{|l|>{\raggedright\arraybackslash}p{4.8cm}|>{\raggedright\arraybackslash}p{4.5cm}|}
\toprule
\textbf{Layanan} & \textbf{Definisi} & \textbf{Mekanisme Kriptografi} \\
\midrule
Confidentiality & Hanya pihak berwenang mengakses & Enkripsi \\
Integrity & Data tidak diubah tanpa izin & Hash, MAC \\
Availability & Data/sistem dapat diakses saat dibutuhkan & Manajemen kunci, redundansi \\
Non-repudiation & Tindakan tidak dapat disangkal & Tanda tangan digital \\
Authenticity & Identitas pengirim/penerima terverifikasi & Sertifikat, protokol auth \\
\bottomrule
\end{tabular}
\caption{Layanan keamanan informasi dan mekanisme kriptografi.}
\end{table}

\begin{contoh}
Dalam pembayaran daring: confidentiality melindungi data kartu kredit; integrity memastikan jumlah transaksi tidak diubah; authenticity memastikan pembayaran dari pemilik kartu yang sah; non-repudiation mencegah penyangkalan transaksi.
\end{contoh}
